\begin{abstract}
A formula is minimal in a logic $L$ if it belongs to $L$ and no strictly
more general formula, with respect to the substitution preorder, belongs to
$L$. We show that a polarized extension-variable translation $\St$ preserves
minimality: if $\alpha$ is minimal in a logic $L$ that contains the axiom
$a\to a$ and the formula $\St(\alpha)$, then $\St(\alpha)$ is also minimal
in $L$. If $L$ is an implicational logic containing $\BCI$, then
$\St(\alpha)\in L$ whenever $\alpha\in L$, and $L\lplus \alpha=L\lplus \St(\alpha)$. Since $\St(\alpha)$
has order at most $4$, we obtain a minimality-preserving normal form: every
implicational logic over $\BCI$ that is axiomatizable by formulas
minimal in classical logic is axiomatizable by such formulas of order at
most $4$, and the bound $4$ is optimal. As a corollary, applying $\St$ to
the counterexample of Nakamura and Matsuda to the Komori--Kashima problem
yields an explicit axiom of Gödel--Dummett logic of order $4$ that is
minimal in classical logic, which answers their Open Problem~2. We also show that the
results extend to extensions of $\FLe$ in
the multiplicative language and of $\FLew$ in the full language.
The main theorem and the soundness of $\St$ over $\BCI$ have been
formalized in Lean~4.
\end{abstract}
